% from ega1-3-fr.tex
\oldpage[I]{112}
l'ensemble produit $P\times Q$ formé des couples $(p,q)$ tels que
$\varphi(p)=\psi(q)$~; on le note $P\times_R Q$. La définition
(\hyperref[I.3.2.1-fr]{3.2.1}) du produit de $S$-préschémas peut encore
s'interpréter, avec les notations de (\hyperref[I.3.4.1-fr]{3.4.1}), par la
formule
\phantomsection\label{I.3.4.2.1-fr}
\[
  (X\times_S Y)(T)=X(T)\times_{S(T)}Y(T)
  \tag{3.4.2.1}
\]
les applications $X(T)\to S(T)$ et $Y(T)\to S(T)$ correspondant aux morphismes
structuraux $X\to S$ et $Y\to S$.
\end{env}

\begin{env}[3.4.3]
\label{I.3.4.3-fr}
Si l'on se donne un préschéma $S$ et que l'on ne considère que les
$S$-préschémas et les $S$-morphismes, on notera encore $X(T)_S$ l'ensemble
$\operatorname{Hom}_S(T,X)$ des $S$-morphismes $T\to X$, en se permettant de
supprimer l'indice $S$ lorsque aucune confusion n'est possible~; on dit encore
que les éléments de $X(T)_S$ sont les \emph{points} (ou \emph{$S$-points}
lorsqu'il y a à craindre des confusions) \emph{du $S$-préschéma $X$ à valeurs
dans le $S$-préschéma $T$}. En particulier, une \emph{$S$-section} de $X$
n'est autre qu'un \emph{point de $X$ à valeurs dans $S$}. La formule
(\hyperref[I.3.4.2.1-fr]{3.4.2.1}) s'écrit alors aussi
\phantomsection\label{I.3.4.3.1-fr}
\[
  (X\times_S Y)(T)_S=X(T)_S\times Y(T)_S~;
  \tag{3.4.3.1}
\]
plus généralement, si $Z$ est un $S$-préschéma, $X$, $Y$, $T$ des
$Z$-préschémas (donc \emph{ipso facto} des $S$-préschémas), on a
\phantomsection\label{I.3.4.3.2-fr}
\[
  (X\times_Z Y)(T)_S=X(T)_S\times_{Z(T)_S}Y(T)_S.
  \tag{3.4.3.2}
\]

On notera que pour démontrer qu'un triplet $(W,r,s)$ formé d'un $S$-préschéma
$W$ et de deux $S$-morphismes $r:W\to X$, $s:W\to Y$ est un produit de $X$ et
$Y$ (au-dessus de $Z$), il suffit par définition de vérifier que pour
\emph{tout $S$-préschéma $T$}, le diagramme
\phantomsection\label{I.3.4.3.product-diagram-fr}
\[
  \xymatrix{
    W(T)_S\ar[r]^{r'}\ar[d]_{s'} & X(T)_S\ar[d]^{\varphi'}\\
    Y(T)_S\ar[r]_{\psi'} & Z(T)_S
  }
\]
fait de $W(T)_S$ le produit fibré de $X(T)_S$ et $Y(T)_S$ au-dessus de
$Z(T)_S$, $r'$ et $s'$ correspondant à $r$ et $s$, $\varphi'$ et $\psi'$ aux
morphismes structuraux $\varphi:X\to Z$, $\psi:Y\to Z$.
\end{env}

\begin{env}[3.4.4]
\label{I.3.4.4-fr}
Lorsque dans ce qui précède, $T$ (resp. $S$) est un schéma affine d'anneau $B$
(resp. $A$), on remplace $T$ (resp. $S$) par $B$ (resp. $A$) dans les notations
précédentes, et on parle alors de \emph{points de $X$ à valeurs dans l'anneau
$B$}, ou de \emph{points du $A$-préschéma $X$ à valeurs dans la $A$-algèbre
$B$} pour les éléments de $X(B)$ et de $X(B)_A$ respectivement. On notera que
$X(B)$ et $X(B)_A$ sont maintenant des foncteurs \emph{covariants en $B$}. On
écrira de même $X(T)_A$ pour l'ensemble des points du $A$-préschéma $X$ à
valeurs dans le $A$-préschéma $T$.
\end{env}

\begin{env}[3.4.5]
\label{I.3.4.5-fr}
Considérons en particulier le cas où $T$ est de la forme $\operatorname{Spec}(A)$,
où $A$ est un anneau \emph{local}~; les éléments de $X(A)$ correspondent alors
biunivoquement aux homomorphismes \emph{locaux} $\mathcal{O}_x\to A$ pour
$x\in X$ (\hyperref[I.2.4.4-fr]{2.4.4})~; on dit que le point $x$ de l'espace
sous-jacent à $X$ est la \emph{localité} du point de $X$ à valeurs dans $A$
auquel il correspond.

Plus particulièrement, on appellera \emph{points géométriques} d'un préschéma
$X$ les \emph{points de $X$ à valeurs dans un corps $K$}~: la donnée d'un tel
point revient donc à la donnée de sa
